\section{Results}

\subsection{First-MU experience states and formal support use}

\Cref{tab:states} shows the primary five-state classification. The largest group is lower strain ($N=3{,}679$), followed by adverse/non-numeric outcomes ($N=1{,}101$), contextual challenge ($N=1{,}012$), individual strain ($N=673$), and compounded strain ($N=162$).

The relationship between academic strain and formal support use is not monotonic. Contextual-challenge students, who are in low-pass-rate classrooms but remain above the individual-strain cutoff, have the highest first-period CMAT-use rate at 28.6\%. Their eventual MU pass rate is 97.7\%, and 76.3\% are observed later in Calculus. By contrast, only 9.5\% of individual-strain students and 8.0\% of compounded-strain students use CMAT in the first MU period, despite much weaker eventual MU pass rates and markedly lower later Calculus progression.

The difference is not explained by the primary cutoff. Across all nine pre-specified combinations of performance and contextual-difficulty thresholds, contextual challenge remains the numeric state with the highest CMAT-use rate. Relative to lower strain, the difference ranges from 8.4 to 10.0 percentage points; relative to individual strain, from 15.7 to 20.1 points; and relative to compounded strain, from 14.5 to 21.8 points.

The threshold-free model gives the same qualitative pattern. Among 5,526 students with numeric relative performance, a one-standard-deviation increase in classroom-relative performance is associated with a 5.0 percentage-point increase in same-period CMAT use ($p<.001$), while a one-standard-deviation increase in classroom pass-rate context is associated with a 4.7-point decrease ($p<.001$). The interaction is negative by 1.8 points ($p<.001$). These estimates describe contemporaneous associations: stronger relative performers are more likely to use formal support, especially in lower-pass-rate environments, but the model cannot establish whether CMAT use preceded, followed, or co-evolved with the realized course outcome.

\subsection{Does a difficult first experience predict later enrolment with a higher-outcome instructor?}

The later Calculus transition does not support the simple hypothesis that students who experience difficulty in MU subsequently compensate by enrolling with instructors who have historically higher outcomes. In the rankable sample of 3,224 later Calculus enrolments, individual-strain students enrol with instructors whose historical-outcome rank is approximately 4.4 percentile points lower than that of lower-strain students ($p=.024$), after adjustment for prior CMAT-use group, degree programme, and Calculus period. The contextual-challenge difference is about $-2.2$ points and is not statistically distinguishable from zero in the primary support specification, while compounded-strain and adverse/non-numeric differences are also imprecise.

This result is stable when the observed period-level instructor set is required to contain more rankable alternatives. Requiring at least two, three, four, or five rankable instructors leaves the individual-strain coefficient near $-4.4$ percentile points, with $p$ values around .024. Under a stricter sample requiring at least 75\% historical instructor coverage, the individual-strain difference is $-6.0$ points ($p=.005$), and contextual challenge is also associated with a lower historical-outcome rank by $-3.4$ points ($p=.019$).

These coefficients should not be interpreted as deliberate preference for a ``harder'' professor. The data identify the instructor with whom the student subsequently enrolled, not the complete set of feasible instructors. The appropriate conclusion is narrower: difficult MU experiences are not followed by a general shift toward historically higher-outcome Calculus instructors in the observed administrative data.

\subsection{Failure as a transition point}

A different pattern appears when the next observation occurs immediately after a failed MU attempt. \Cref{tab:failure} summarizes repeated-attempt transitions. After the first failed or adverse attempt, 97.0\% of students change instructor. Among the 387 transitions for which the previous and next instructors can both be placed on a comparable historical-outcome percentile, 62.3\% move upward, with an average increase of 13.4 percentile points. The same direction appears after the second and third failed attempts, although the sample becomes much smaller.

The movement is not an artefact of the within-period percentile denominator. After the first failed attempt, 67.5\% of 456 comparable transitions move to an instructor with a higher strictly-prior pass rate, while 66.9\% of 387 move to an instructor with a higher strictly-prior mean grade. After the second failure, the corresponding shares are 62.4\% and 58.8\%; after the third, 65.4\% and 68.2\%. This consistency across absolute and relative historical measures supports the descriptive conclusion that failed students commonly re-enter MU in an instructor context associated with stronger prior academic outcomes.

CMAT use follows a different pattern. After the first failure, any CMAT use falls from 18.9\% in the failed-attempt period to 9.8\% in the next attempt, and the mean number of visits falls by 0.23. Attempt number itself is not jointly associated with the three modeled adaptation outcomes after adjustment. Prior CMAT use, however, is associated with a 11.3 percentage-point higher probability of CMAT use in the next attempt in the full transition sample ($p<.001$), although it does not predict an increase in the number of visits.

Among the subset of numeric failures with a valid classroom-relative $Z$, failure severity provides little evidence of a simple dose-response adaptation mechanism. A one-standard-deviation increase in prior $Z$ is associated with a 5.3 percentage-point higher probability of moving upward in historical instructor context, a borderline association ($p=.056$), while its associations with next-attempt CMAT use and increased CMAT use are small and statistically non-significant. Thus, the descriptive movement after failure should not be interpreted as ``more severe failure produces more adaptation.''

\subsection{Degree-programme heterogeneity}

Degree programme contributes additional information to each major behavioural outcome, but the incremental explanatory magnitude is modest; \Cref{tab:degree} summarizes the three omnibus comparisons. For first-MU CMAT use, adding degree programme to experience state and period increases $R^2$ by 0.0154, with a joint programme test of $p<.001$. For later Calculus CMAT use, the corresponding increase is 0.0250 ($p<.001$), and for the historical outcome rank of the subsequently enrolled Calculus instructor it is 0.0143 ($p<.001$).

The programme-by-experience-state interaction is also detectable when smaller programmes are pooled before estimation. With programmes retained separately only when they contain at least 100 students, the interaction adds $\Delta R^2=0.0112$; at thresholds of 150 and 200 students, the increments fall to 0.0087 and 0.0060, respectively, while the joint interaction tests remain statistically clear. This pattern is consistent with disciplinary heterogeneity in behavioural response but does not suggest that degree programme dominates the individual and contextual dimensions of the first-MU experience.

Actuar\'ia provides a useful descriptive example. Among 368 first-MU students in the programme, 30.4\% use CMAT in the first MU period and 40.8\% of the 287 students later observed in Calculus use CMAT there. Yet the mean historical-outcome percentile of their Calculus instructor is 0.532, close to the middle of the observed distribution. High formal-support use in this programme is therefore not simply accompanied by systematic enrolment with historically high-outcome Calculus instructors.

\draftnote{Decide at journal-targeting stage whether the Actuar\'ia example belongs in the main text or in supplementary material; the omnibus programme result should remain primary.}

